% --- MAGIC COMMENTS: CONFIGURAZIONE EDITOR --- (vedi main.tex) ---
% !TEX root = ../main.tex
\chapter{Prefazione, ma da leggere!}
\usection{L'obiettivo Principale}
Questo libro (se così può essere chiamato) è rivolto principalmente agli
studenti del \course{} in \emph{\prog} del \emph{\uni}, ed è stato pensato per
aiutare costoro a preparare l'esame di \emph{\DDcoursename}, presente nel
programma del terzo anno al primo semestre.

Nasce in precisione come una sorta di blocco appunti avanzato, da integrare
ovviamente al materiale didattico del corso, del quale necessariamente non può
essere un sostituto completo! Si invitano dunque tutti gli studenti che
intendano farne uso, a \textcolor{Red}{\textbf{consultare sempre e comunque il
materiale ufficiale fornito dai docenti}}, in modo da potersi preparare al
meglio e secondo le linee guida.

\usection{L'obiettivo Secondario}
Fatta la doverosa premessa iniziale è chiaro che, a causa della sua natura,
quest'opera non possa porsi come obiettivo quello di essere un riferimento
completo e definitivo per lo studio dell'elettronica (di base). Si ribadisce
nuovamente che il suo scopo è quello di essere un supporto allo studio, mirato
al superamento dell'esame.

Tuttavia, sebbene questi aspetti puramente utilitaristici siano alla base della
sua esistenza, parlare solo di essi non renderebbe giustizia alla filosofia e
alla visione didattica con cui è stata pensata e realizzata. Infatti, se da un
lato l'elettronica è una materia ormai consolidata, per la quale è possibile
reperire (anche troppo) facilmente buon materiale per lo studio; dall'altro gli
studenti faticano a orientarsi efficacemente in questa vastità, e a sostenere un
approccio che sia produttivo e accademicamente sano allo stesso tempo.

Dall'esperienza dell'autore, è evidente che le difficoltà principali vengano
alla luce quando si devono far convivere assieme:
\begin{itemize}
\item il rigore matematico (e logico) necessario a una comprensione chiara, ed
eventualmente profonda, dei modelli studiati;
\item un approccio \fig{fisico} e intuitivo, per non perdersi nei dettagli e
costruire un motore mentale di \fig{alto livello} (in senso informatico), con
cui navigare gli argomenti e i concetti;
\item le capacità pratiche e le conoscenze tecniche richieste nei laboratori,
nei quali tipicamente è molto raro che le cose funzioni al primo tentativo,
nonostante la preparazione teorica.
\end{itemize}

È ovvio che non possa esistere un singolo corso universitario, né un singolo
libro, che da solo possa insegnare efficacemente tutto questo; cioè
\emph{come studiare} e gli \emph{strumenti per studiare}, oltre a \emph{cosa
studiare} e \emph{come applicare quanto studiato}. Si noti tra l'altro che sono
tutte cose ben diverse, e a volte lontane, dal \emph{superare un esame}.

In ogni caso, è proprio in questa posizione ibrida che la presente guida vuole
tentare di porsi: l'obiettivo secondario, per quanto ambizioso possa essere, è
quello di provare a ridurre almeno di un po' le distanze che esistono tra gli
obiettivi materiali, e una buona preparazione personale. Sarà sicuramente
complesso, ma rimane viva la speranza che possa essere un aiuto diverso e più
vicino alle esigenze reali degli studenti, che vogliono sì superare l'esame, ma
anche \emph{capire meglio}, o \emph{capire perché non capiscono}.

\usection{Stile e Contenuti del Testo}
Essendo stato pensato più come una guida allo studio che come un testo classico,
lo stile che si è tentato di adottare è asciutto e diretto, per favorire la
comprensione e la sintesi, ma non per questo banale o superficiale. Chi vorrà,
troverà spazio per maggiori dettagli e approfondimenti utili, opportunamente
segnalati come argomenti a parte, o con annotazioni specifiche.

Il contenuto del libro cerca inoltre di \textbf{rispecchiare fedelmente il
contenuto e il programma del corso}, tant'è che i suoi capitoli sono stati
denominati e suddivisi seguendo lo stesso schema utilizzato dal materiale
didattico \autocite{DispenseElettronica}, in precisione dalle slide (o per i più
nostalgici, \fig{dai lucidi}). Si confida nel fatto che questo possa favorire
una sinergia generale con esso, e quindi che possa aiutare lo studente a
orientarsi e a trarne benefici.

Per tali ragioni, ogni argomento trattato è da considerarsi come \textbf{parte
effettiva del programma, tranne dove diversamente indicato}, ad esempio in caso
delle digressioni o degli approfondimenti citati.

Le altre principali risorse che sono state sfruttate per la stesura, sono le
stesse indicate dai docenti come riferimenti esterni ufficiali, in particolare:
\begin{itemize}
\item il libro di testo di \emph{\en{Storey}}
\autocite{storey4ed}\autocite{storey6ed} (quello principale)
\item il libro di testo di \emph{\en{Floyd}} \autocite{floyd}, e quello degli
autori \emph{\en{Jaeger, Blalock e Blalock}} \autocite{jaegerblalock} (quelli
secondari)
\end{itemize}

Eventuali ulteriori riferimenti bibliografici (anche aggiunti successivamente al
momento della scrittura della presente prefazione), saranno opportunamente
riportati nella \textbf{\nameref{chap:bibliography}} presente alla fine.

\usection{Sì, ma quanta Matematica?}
Infine, una nota di incoraggiamento per tutti coloro che possano averne bisogno:
è stato citato il \quo{rigore matematico}, un concetto spesso frainteso e
ritenuto più antipatico che utile. Si vuole ora rincuorare lo studente: qui se
ne intende una \emph{versione sana}, che con il giusto sforzo iniziale
(inevitabile), permette poi di comprendere serenamente e facilmente, ma
soprattutto una volta per tutte, ognuno dei concetti affrontati. Non ci si lasci
dunque spaventare dalle equazioni, dai teoremi o dalle dimostrazioni, perché non
c'è molto da temere (se sono state spiegate chiaramente).

Nello studio delle materie scientifiche, la matematica è una forma di
investimento che diventa conveniente molto rapidamente con il tempo: saper
formalizzare con i simboli giusti, vale molto più di quanto tante parole e
\fig{reinterpretazioni pratiche} possano dire al loro posto: i simboli portano
con sé un significato \fig{completo}, una volta che sono stati \emph{ben
definiti}. E a volte anche al di là delle aspettative%
\footnote{Il fisico \en{Eugene Wigner} parlava letteralmente de
\emph{L'irragionevole efficacia della matematica} \autocite{wigner}.}.
Non a caso i maggiori progressi nella storia della conoscenza umana, ci sono
stati quando si è passati da un approccio \emph{retorico}%
\footnote{Matematica Retorica, ossia \lit{parlata}: sviluppata principalmente
attraverso il linguaggio naturale, inclusi calcoli e dimostrazioni. Utilizzata
dagli antichi Greci e Arabi, sopravvissuta fino al Medioevo in Europa. Il
passaggio alla Matematica Simbolica è avvenuto solo dopo il 1500, grazie a
figure come \fr{Viète} e Cartesio. Il passaggio ovviamente non è stato netto, e
la forma di intermedia di transizione è stata chiamata \emph{Matematica
Sincopata}.} a uno \emph{simbolico}.
Quindi, perché tornare indietro?

Per fare un paragone, che potrebbe essere utile ad apprezzare lo sforzo di chi
ha scritto il presente documento, è un po' lo stesso fenomeno che si vive
quando, scrivendo un documento complesso, si passa dall'utilizzare
\emph{\en{Word}}, o qualsiasi altro editor \emph{\en{WYSIWYG}}%
\footnote{Sta per \emph{\en{What You See Is What You Get}}: costringe
l'utente a occuparsi degli aspetti grafici assieme al contenuto.},
a utilizzare \LaTeX, che è un motore tipografico di tipo \emph{\en{WYSIWYM}}%
\footnote{Sta per \emph{\en{What You See Is What You Mean}}: gli aspetti grafici
sono gestiti in automatico, ma serve conoscere del codice da aggiungere al
testo.}.
La curva di apprendimento iniziale è piuttosto ripida, ma una volta
superata, il risparmio di energie e risorse è molto grande. Si veda
\cref{fig:wordVSlatex.png}. \addfigure[Word vs \LaTeX]{wordVSlatex.png}{0.4}%

Ovviamente, tutto ciò non toglie che sia necessario sviluppare un approccio
intuitivo in parallelo all'approccio formale, fondamentale per la padronanza.

% Da aggiungere: origine immagine. O rifai immagine.
